\begin{lemma}[Left presentation invariance]
If $\mathsf{PowerPresentationEq}(X,X')$, then for every
$Y\in\mathsf{PowerQ}(Q)$,
\[
\mathsf{PowerRel}(r)(X,Y)\quad\Longleftrightarrow\quad
\mathsf{PowerRel}(r)(X',Y).
\]
\end{lemma}
\begin{proof}
Induct on the well-founded child descent of the compared left
presentations, using the recursive definition in \noderef{PowerRel}.  If
both left presentations are atoms, \noderef{PowerPresentationEq} gives
equality of their $Q$-values, so the atom--atom and atom--node reductions
from \noderef{PowerRelAtomAtom} and \noderef{PowerRelAtomNode} give the same
proposition.  A mixed atom--node pair cannot satisfy the equivalence
hypothesis, since that branch of \noderef{PowerPresentationEq} is false.

In the remaining case both presentations are nodes.  For each child of the
first node, the forward conjunct in \noderef{PowerPresentationEq} supplies a
child of the second node structurally equivalent to it; the induction
hypothesis transfers the recursive comparison with any fixed right child.
This transports the existential right-child witnesses in the node--node
reduction \noderef{PowerRelNodeNode}.  The node--atom reduction
\noderef{PowerRelNodeAtom} is handled in the same way, with the universal
quantifier over left children.  The reverse implication uses the backward
matching conjunct.  Symmetry and transitivity of structural matching are
available from \noderef{PowerPresentationEqEquivalence}, so the induction
remains valid when matching witnesses are composed.
\end{proof}
